% GENERATED FILE. Edit canonical lessons, then run npm run generate:books.
% canonical-source-hash: fnv1a64:4461e809195e3990
% canonical-lessons: AR-C77-najm, AR-C77-sahaba, AR-C77-shams, AR-C77-qamar, AR-C77-bahr

\chapter{Star, Cloud, Sun, Moon, Sea}
\label{ch:ar-star-cloud-sun-moon-sea}
\hlchaptermodality{77}

\begin{chapteropening}
\textbf{By the end of this chapter:} \emph{I can say a star, a cloud, the sun, the moon and the sea.}

Complete the last lesson of chapter 77: I can say a star, a cloud, the sun, the moon and the sea.
\end{chapteropening}

\section[\texorpdfstring{\ar{نجم} (najm)}{نجم (najm)}]{\ar{نجم} (najm) — a star}

\label{lesson:AR-C77-najm}



\begin{quote}
Before the new one: say the Arabic for sand, then the Arabic for a field.
\end{quote}

\subsection*{You'll want to know: \ar{نجم}}
\textbf{\ar{نجم}} — \emph{najm} — \textquotedblleft{}a star\textquotedblright{}.

\begin{cousinweb}
\textbf{\ar{نجم}} (\emph{najm}) is built on the native Arabic root \emph{n-j-m}, \textquotedblleft{}a star; to rise, to appear\textquotedblright{} (Lane). Lane gives \emph{najm} as a star, and also the Pleiades.
\end{cousinweb}

\begin{grammarlens}[title={What you've built}]
One more word for this chapter.
\end{grammarlens}

\subsection*{Your turn}
\begin{itemize}
\raggedright
  \item \emph{Say it:} \emph{najm}
  \item \emph{Say it:} \emph{najm}, once more
  \item \emph{Say it:} say \emph{najm}
  \item \emph{Recall:} say the Arabic for sand, then the Arabic for a field, then say \emph{najm} again
\end{itemize}

\begin{tcolorbox}[breakable,colback=teal!4,colframe=teal!35!black,title={Before you move on}]
What does \ar{نجم} mean? (\textquotedblleft{}A star\textquotedblright{}.) Say it once more.
\end{tcolorbox}
\section[\texorpdfstring{\ar{سحابة} (saḥāba)}{سحابة (saḥāba)}]{\ar{سحابة} (saḥāba) — a cloud}

\label{lesson:AR-C77-sahaba}



\begin{quote}
Before the new one: say the Arabic for a field, then the Arabic for a star.
\end{quote}

\subsection*{You'll want to know: \ar{سحابة}}
\textbf{\ar{سحابة}} — \emph{saḥāba} — \textquotedblleft{}a cloud\textquotedblright{}.

\begin{cousinweb}
\textbf{\ar{سحابة}} (\emph{saḥāba}) is built on the native Arabic root \emph{s-ḥ-b}, \textquotedblleft{}to drag, to draw along\textquotedblright{} (Lane). \emph{Saḥāba} is one cloud of \emph{saḥāb}, \textquotedblleft{}clouds\textquotedblright{}.
\end{cousinweb}

\begin{grammarlens}[title={What you've built}]
One more word for this chapter.
\end{grammarlens}

\subsection*{Your turn}
\begin{itemize}
\raggedright
  \item \emph{Say it:} \emph{saḥāba}
  \item \emph{Say it:} \emph{saḥāba}, once more
  \item \emph{Say it:} say \emph{saḥāba}
  \item \emph{Recall:} say the Arabic for a field, then the Arabic for a star, then say \emph{saḥāba} again
\end{itemize}

\begin{tcolorbox}[breakable,colback=teal!4,colframe=teal!35!black,title={Before you move on}]
What does \ar{سحابة} mean? (\textquotedblleft{}A cloud\textquotedblright{}.) Say it once more.
\end{tcolorbox}
\section[\texorpdfstring{\ar{شمس} (shams)}{شمس (shams)}]{\ar{شمس} (shams) — the sun}

\label{lesson:AR-C77-shams}



\begin{quote}
Before the new one: say the Arabic for a star, then the Arabic for a cloud.
\end{quote}

\subsection*{You'll want to know: \ar{شمس}}
\textbf{\ar{شمس}} — \emph{shams} — \textquotedblleft{}the sun\textquotedblright{}.

\begin{cousinweb}
\textbf{\ar{شمس}} (\emph{shams}) is built on the native Arabic root \emph{sh-m-s}, \textquotedblleft{}the sun\textquotedblright{} (Lane). Hebrew \emph{shemesh}, \textquotedblleft{}sun\textquotedblright{}, is its cousin.
\end{cousinweb}

\begin{grammarlens}[title={What you've built}]
One more word for this chapter.
\end{grammarlens}

\subsection*{Your turn}
\begin{itemize}
\raggedright
  \item \emph{Say it:} \emph{shams}
  \item \emph{Say it:} \emph{shams}, once more
  \item \emph{Say it:} say \emph{shams}
  \item \emph{Recall:} say the Arabic for a star, then the Arabic for a cloud, then say \emph{shams} again
\end{itemize}

\begin{tcolorbox}[breakable,colback=teal!4,colframe=teal!35!black,title={Before you move on}]
What does \ar{شمس} mean? (\textquotedblleft{}The sun\textquotedblright{}.) Say it once more.
\end{tcolorbox}
\section[\texorpdfstring{\ar{قمر} (qamar)}{قمر (qamar)}]{\ar{قمر} (qamar) — the moon}

\label{lesson:AR-C77-qamar}



\begin{quote}
Before the new one: say the Arabic for a cloud, then the Arabic for the sun.
\end{quote}

\subsection*{You'll want to know: \ar{قمر}}
\textbf{\ar{قمر}} — \emph{qamar} — \textquotedblleft{}the moon\textquotedblright{}.

\begin{cousinweb}
\textbf{\ar{قمر}} (\emph{qamar}) is built on the native Arabic root \emph{q-m-r}, \textquotedblleft{}the moon\textquotedblright{} (Lane). An old noun with its own root in Lane.
\end{cousinweb}

\begin{grammarlens}[title={What you've built}]
One more word for this chapter.
\end{grammarlens}

\subsection*{Your turn}
\begin{itemize}
\raggedright
  \item \emph{Say it:} \emph{qamar}
  \item \emph{Say it:} \emph{qamar}, once more
  \item \emph{Say it:} say \emph{qamar}
  \item \emph{Recall:} say the Arabic for a cloud, then the Arabic for the sun, then say \emph{qamar} again
\end{itemize}

\begin{tcolorbox}[breakable,colback=teal!4,colframe=teal!35!black,title={Before you move on}]
What does \ar{قمر} mean? (\textquotedblleft{}The moon\textquotedblright{}.) Say it once more.
\end{tcolorbox}
\section[\texorpdfstring{\ar{بحر} (baḥr)}{بحر (baḥr)}]{\ar{بحر} (baḥr) — the sea}

\label{lesson:AR-C77-bahr}



\begin{quote}
Before the new one: say the Arabic for the sun, then the Arabic for the moon.
\end{quote}

\subsection*{You'll want to know: \ar{بحر}}
\textbf{\ar{بحر}} — \emph{baḥr} — \textquotedblleft{}the sea\textquotedblright{}.

\begin{cousinweb}
\textbf{\ar{بحر}} (\emph{baḥr}) is built on the native Arabic root \emph{b-ḥ-r}, \textquotedblleft{}the sea; to slit\textquotedblright{} (Lane). Lane also gives a verb \emph{baḥara}, \textquotedblleft{}to slit, to cleave\textquotedblright{}.
\end{cousinweb}

\begin{grammarlens}[title={What you've built}]
One more word for this chapter.
\end{grammarlens}

\subsection*{Your turn}
\begin{itemize}
\raggedright
  \item \emph{Say it:} \emph{baḥr}
  \item \emph{Say it:} \emph{baḥr}, once more
  \item \emph{Say it:} say \emph{baḥr}
  \item \emph{Recall:} say the Arabic for the sun, then the Arabic for the moon, then say \emph{baḥr} again
\end{itemize}

\begin{tcolorbox}[breakable,colback=teal!4,colframe=teal!35!black,title={Before you move on}]
What does \ar{بحر} mean? (\textquotedblleft{}The sea\textquotedblright{}.) Say it once more.
\end{tcolorbox}
